% -*- coding: utf-8 -*-
% !TEX program = xelatex
\documentclass[plain,noanswer,medmath]{../../template/jnuexam}
\usepackage{../../template/kysx}

\begin{document}


\examtitle{2015}{1}


\makepart{选择题}{1～8小题，每小题4分，共32分}


\begin{problem}
设函数$f(x)$在$\left(-\infty , +\infty \right)$内连续，
其中二阶导数$f''(x)$的图形如图所示，则曲线$y = f(x)$的拐点的个数为\pickout{}
\par\centerline{\includegraphics[width=0.5\linewidth]{2015A}}
\begin{abcd}
\item $0$．
\item $1$．
\item $2$．
\item $3$．
\end{abcd}
\end{problem}


\begin{problem}
设$y = \frac{1}{2}\e^{2x} + \left(x - \frac{1}{3} \right) \e^x$
是二阶常系数非齐次线性微分方程$y'' + ay' + by = c\e^x$的一个特解，则\pickout{}
\begin{abcd}
\item $a = - 3,b = 2,c = - 1$．
\item $a = 3,b = 2,c = - 1$．
\item $a = - 3,b = 2,c = 1$．
\item $a = 3,b = 2,c = 1$．
\end{abcd}
\end{problem}


\begin{problem}
若级数$\sum_{n = 1}^{\infty}a_n$条件收敛，则$x = \sqrt{3}$与$x = 3$
依次为幂级数$\sum_{n = 1}^{\infty}na_n (x - 1)^n$的\pickout{}
\begin{abcd}
\item 收敛点，收敛点．
\item 收敛点，发散点．
\item 发散点，收敛点．
\item 发散点，发散点．
\end{abcd}
\end{problem}


\begin{problem}
设$D$是第一象限由曲线$2xy = 1$，$4xy = 1$，与直线$y = x$，$y = \sqrt{3} x$围成的平面区域，
函数$f(x,y)$在$D$上连续，则$\iint_D f(x,y)\dx\dy =$\pickout{}
\begin{abcd}
\item $\int_{\frac{\pi}{4}}^{\frac{\pi}{3}} \d\theta
       \int_{\frac{1}{2\sin 2\theta}}^{\frac{1}{\sin 2\theta}} f(r\cos \theta ,r\sin \theta)r\dr$．
\item $\int_{\frac{\pi}{4}}^{\frac{\pi}{3}} \d\theta
       \int_{\frac{1}{\sqrt{2\sin 2\theta}}}^{\frac{1}{\sqrt{\sin 2\theta}}}
            f(r\cos \theta ,r\sin \theta)r\dr$．
\item $\int_{\frac{\pi}{4}}^{\frac{\pi}{3}} \d\theta
       \int_{\frac{1}{2\sin 2\theta}}^{\frac{1}{\sin 2\theta}} f(r\cos \theta ,r\sin \theta)\dr$．
\item $\int_{\frac{\pi}{4}}^{\frac{\pi}{3}} \d\theta
       \int_{\frac{1}{\sqrt{2\sin 2\theta}}}^{\frac{1}{\sqrt{\sin 2\theta}}}
            f(r\cos \theta ,r\sin \theta)\dr$．
\end{abcd}
\end{problem}


\begin{problem}
设矩阵$A = \begin{pmatrix}
  1 & 1 & 1 \\
  1 & 2 & a \\
  1 & 4 & a^2
\end{pmatrix}$，$b = \begin{pmatrix}
  1 \\
  d \\
  d^2
\end{pmatrix}$．若集合$\Omega = \left\{1,\left. 2 \right\} \right.$，
则线性方程组$Ax = b$有无穷多解的充分必要条件为\pickout{}
\begin{abcd}
\item $a \notin \Omega ,d \notin \Omega$．
\item $a \notin \Omega ,d \in \Omega$．
\item $a \in \Omega ,d \notin \Omega$．
\item $a \in \Omega ,d \in \Omega$．
\end{abcd}
\end{problem}


\begin{problem}
设二次型$f\left(x_1,x_2,x_3 \right)$在正交变换为$x = Py$下的标准形为$2y_1^2 + y_2^2 - y_3^2$，
其中$P = \left(e_1,e_2,e_3 \right)$，若$Q = \left(e_1, - e_3,e_2 \right)$，
则$f\left(x_1,x_2,x_3 \right)$在正交变换$x = Qy$下的标准形为\pickout{}
\begin{abcd}
\item $2y_1^2 - y_2^2 + y_3^2$．
\item $2y_1^2 + y_2^2 - y_3^2$．
\item $2y_1^2 - y_2^2 - y_3^2$．
\item $2y_1^2 + y_2^2 + y_3^2$．
\end{abcd}
\end{problem}


\begin{problem}
若$A$, $B$为任意两个随机事件，则\pickout{}
\begin{abcd}
\item $P(AB) \le P(A)P(B)$．
\item $P(AB) \ge P(A)P(B)$．
\item $P(AB) \le \frac{P(A) + P(B)}{2}$．
\item $P(AB) \ge \frac{P(A) + P(B)}{2}$．
\end{abcd}
\end{problem}


\begin{problem}
设随机变量$X,Y$不相关，且$EX = 2$，$EY = 1$，$DX = 3$，
则$E\left[X\left(X + Y - 2 \right) \right] =$\pickout{}
\begin{abcd}
\item $- 3$．
\item $3$．
\item $- 5$．
\item $5$．
\end{abcd}
\end{problem}


\makepart{填空题}{9～14小题，每小题4分，共24分}


\begin{problem}
$\lim_{x \to 0} \frac{\ln \cos x}{x^2} =$ \fillin{}．
\end{problem}


\begin{problem}
$\int_{-\frac{\pi}{2}}^{\frac{\pi}{2}} \left(\frac{\sin x}{1+\cos x} + |x|\right)\dx =$ \fillin{}．
\end{problem}


\begin{problem}
若函数$z = z(x,y)$由方程$e^x + xyz + x + \cos x = 2$确定，则$\dz\big|_{(0,1)} =$ \fillin{}．
\end{problem}


\begin{problem}
设$\Omega$是由平面$x + y + z = 1$与三个坐标平面平面所围成的空间区域，\goodbreak
则$\iiint_{\Omega}(x + 2y + 3z)\dx\dy\dz =$ \fillin{}．
\end{problem}


\begin{problem}
$n$阶行列式$\begin{vmatrix}
       2 &      0 & \cdots &      0 & 2 \\
      -1 &      2 & \cdots &      0 & 2 \\
  \vdots & \vdots & \ddots & \vdots & \vdots \\
       0 &      0 & \cdots &      2 & 2 \\
       0 &      0 & \cdots &     -1 & 2
\end{vmatrix} =$ \fillin{}．
\end{problem}


\begin{problem}
设二维随机变量$(X,Y)$服从正态分布$N(1,0;1,1,0)$，则$P\{XY - Y < 0\} =$ \fillin{}．
\end{problem}


\makepart{解答题}{15～23题，共94分}


\begin{problem}[points=10]
设函数$f(x) = x + a\ln(1 + x) + bx\sin x$，$g(x) = kx^3$，
若$f(x)$与$g(x)$在$x \to 0$是等价无穷小，求$a,b,k$的值．
\end{problem}


\begin{problem}[points=10]
设函数$f(x)$在定义域$I$上的导数大于零，若对任意的$x_0 \in I$，
由线$y = f(x)$在点$\left(x_0,f(x_0)\right)$处的切线与直线$x = x_0$及$x$轴所围成区域的面积恒为$4$，
且$f(0) = 2$，求$f(x)$的表达式．
\end{problem}


\begin{problem}[points=10]
已知函数$f(x,y) = x + y + xy$，曲线$C: x^2 + y^2 + xy = 3$，
求$f(x,y)$在曲线C上的最大方向导数．
\end{problem}


\begin{problem}[points=10]
\begin{enumerate}
\item
设函数$u(x),v(x)$可导，利用导数定义证明$[u(x)v(x)]' = u'(x)v(x) + u(x)v'(x)$．
\item
设函数$u_1(x),u_2(x), \cdots ,u_n(x)$可导，$f(x) = u_1(x) u_2(x) \cdots u_n(x)$，写出$f(x)$的求导公式．
\end{enumerate}
\end{problem}


\begin{problem}[points=10]
已知曲线$L$的方程为$\begin{cases}
  z = \sqrt{2 - x^2 - y^2} , \\
  z = x,
\end{cases}$起点为$A\left(0,\sqrt{2} ,0 \right)$，终点为$B\left(0, - \sqrt{2} ,0 \right)$，
计算曲线积分$I = \int_L (y + z)\dx + \left(z^2 - x^2 + y \right)\dy + (x^2 + y^2)\dz$．
\end{problem}


\begin{problem}
（本题满11分）
设向量组$\alpha_1,\alpha_2,\alpha_3$为$R^3$的一个基，
$$\beta_1 = 2\alpha_1 + 2k\alpha_3,\quad \beta_2 = 2\alpha_2,\quad
  \beta_3 = \alpha_1 + (k + 1)\alpha_3.$$
\begin{enumerate}
\item
证明向量组$\beta_1,\beta_2,\beta_3$为$R^3$的一个基；
\item
当$k$为何值时，存在非零向量$\xi$在基$\alpha_1,\alpha_2,\alpha_3$
与基$\beta_1,\beta_2,\beta_3$下的坐标相同，并求所有的$\xi$．
\end{enumerate}
\end{problem}


\begin{problem}[points=11]
设矩阵$A = \begin{pmatrix}
   0 &  2 & -3 \\
  -1 &  3 & -3 \\
   1 & -2 & a
\end{pmatrix}$相似于矩阵$B = \begin{pmatrix}
  1 & -2 & 0 \\
  0 &  b & 0 \\
  0 &  3 & 1
\end{pmatrix}$．\par
\begin{enumerate*}
\item 求$a,b$的值；
\item 求可逆矩阵$P$，使$P^{-1} AP$为对角矩阵．
\end{enumerate*}
\end{problem}


\begin{problem}[points=11]
设随机变量$X$的概率密度为$f(x) = \begin{cases}
  2^{- x}\ln 2, & x > 0, \\
             0, & x \le 0.
\end{cases}$对$X$进行独立重复的观测，直到第$2$个大于$3$的观测值出现时停止．记$Y$为观测次数．\par
\begin{enumerate*}
\item 求$Y$的概率分布；
\item 求$EY$．
\end{enumerate*}
\end{problem}


\begin{problem}[points=11]
设总体X的概率密度为
$$f(x,\theta) = \begin{cases}
  \frac{1}{1 - \theta}, & \theta \le x \le 1, \\
                     0, & \text{其他}.
\end{cases}$$
其中$\theta$为未知参数，$x_1,x_2, \cdots ,x_n$为来自该总体的简单随机样本．\par
\begin{enumerate*}
\item 求$\theta$的矩估计量；
\item 求$\theta$的最大似然估计量．
\end{enumerate*}
\end{problem}


\end{document}
